\chapter{Conclusions}\label{ch:conclusions}
This thesis addressed the filtering and categorization of cybersecurity-related social media posts within the CS-Aware-Next project. Two architectures were built and compared on the same data: a Two-Step System that first filters the relevant posts and then categorizes them, and a One-Step System that categorizes all the posts with a single top-level model and implies relevance from the predicted categories. Both were developed on one grouped train-test split, with selection and tuning confined to the training partition, and evaluated once on the test partition with 95{\%} group-bootstrap confidence intervals.

The filter reaches a balanced accuracy of 0.787 (95{\%} interval from 0.769 to 0.808) and a ROC AUC of 0.876 (from 0.859 to 0.891) on the test set. The embedder of each component was selected by cross-validation, and the best one varies by component, with small differences among the candidates in the selection scores, so the choice should not be read as a ranking of the embedders. The categorization at the sub levels shows lower and wider-interval scores than the filter, and the labels there have a small support.

On the end-to-end comparison on the public labels, the micro F1 of the Two-Step System (0.651) and of the One-Step System (0.659) are not distinguishable: the paired difference is $-$0.007 with a 95{\%} interval from $-$0.030 to 0.011. The systems differ in the balance between precision and recall. The One-Step System has the higher micro precision (0.694 against 0.644), and the Two-Step System the higher micro recall (0.659 against 0.627). The recall of the Two-Step System is bounded by the recall of its filter, 0.798, since the relevant posts rejected by the filter lose their mapped labels. Which system is preferable therefore depends on the cost of a missed post relative to the cost of a false alert, and no claim of superiority is supported. The subset accuracy gives weak evidence in favour of the One-Step System ($-$0.020, interval from $-$0.042 to $-$0.001, without adjustment for multiple comparisons). Efficiency was not measured, and nothing is claimed about it. The evaluation on NistPosts, a dataset with a different annotation, is descriptive: both systems obtain high precision, and the One-Step System has the lower recall.

Several limitations qualify these findings. The results come from a single split, so they do not capture the variability due to the choice of the partition. The confidence intervals are conditional on the frozen trained systems and reflect only the resampling of the test posts. The category \texttt{Other} covers both the posts without a relevant category and the relevant posts without a mapped category, which affects the implicit relevance of the One-Step System and the interpretation of the public labels. The labels of the sub levels have a small support, so their estimates are wide. NistPosts follows a different annotation and cannot be used to compare the two systems on equal terms.

These limits suggest the directions for future work. Repeating the evaluation over several grouped splits, or nested resampling, would quantify the variability of the partition. A taxonomy that separates irrelevant posts from relevant posts without a mapped category would make the semantics of \texttt{Other} explicit. Additional annotated posts for the rare sub-level labels, or techniques designed for rare labels such as cost-sensitive learning, could improve their estimates. A decision threshold chosen for the desired balance of precision and recall would allow the operating point to be matched to an organization's needs, and a dedicated measurement of the computational cost would allow the two architectures to be compared on efficiency. Finally, annotating a second dataset with the same taxonomy would support a controlled evaluation of the generalization of both systems.
